\section{Introduction}
\label{sec:intro}

The processes studied in this thesis describe populations small enough that their
fate is decided by chance rather than by their mean behaviour --- a founding cohort
of a few individuals, or the contents of a single compartment that may later
rupture. At those counts a deterministic rate equation answers the wrong question.
It returns a single trajectory, growing or decaying, when what one wants to know is
whether the population establishes itself at all, and how large it is if it does.
Both are questions about a distribution, and neither is visible in a mean.

Two threads run through this chapter. They are worth separating at the outset,
because almost every later section belongs to one or the other.

The first is \emph{conditioning on survival}. A subcritical branching process dies
out with probability one, so its unconditional mean decays geometrically and
records little beyond the certainty of extinction. Restrict attention to those
realisations that happen still to be alive at generation $n$, however, and the
picture changes entirely: the conditional population settles to a fixed
distribution, and the mean of that distribution is a constant. For the binary
Galton--Watson process with replication probability $p<\tfrac12$ the constant is
$1/A(p)$, where
\begin{equation*}
A(p) = \lim_{n\to\infty}\frac{S_n}{(2p)^n}
\end{equation*}
and $S_n$ is the probability of survival to generation $n$. A large part of this
chapter is about $A(p)$: what it is, how to compute it, and why no elementary
formula for it has been found despite a determined search.

The second thread is the \emph{near-critical regime}. As $p$ rises towards
$\tfrac12$ the process takes ever longer to die, the quasi-stationary mean $1/A(p)$
diverges, and any numerical scheme that works by iterating to stationarity becomes
impractical. Just above $\tfrac12$ the size of the founding cohort begins to weigh
heavily on whether the process survives at all. The behaviour on either side of
$p=\tfrac12$ is what makes small founding populations interesting, and it is also
what makes them awkward to compute with.

\Cref{sec:prelim} assembles the standard apparatus: exponential clocks and the
races between them, continuous-time Markov chains and their generators,
probability generating functions, and the linear birth--death process, whose
generating-function partial differential equation is the template for everything in
\cref{sec:moc}. \Cref{sec:gw} sets out the Galton--Watson process and establishes
the two facts on which the rest depends --- that extinction is certain at and below
criticality, and that at criticality the expected lifetime is nevertheless infinite.
\Cref{sec:qsd} derives the limiting conditional mean and variance, and so
introduces $A(p)$. \Cref{sec:small} asks what changes when the process starts from a
cohort of $k$ individuals rather than one, treats a discrete birth--death process
subject to catastrophe, obtains the corresponding constant for the continuous-time
birth--death process --- where a closed form \emph{is} available --- and defines the
continuous-time birth--death--catastrophe process used in later chapters.
\Cref{sec:Ap} is devoted to $A(p)$ itself. \Cref{sec:moc} then changes tack and
develops the method of characteristics for the generating-function equations of
absorption models.

Two remarks on how the chapter is written. First, \cref{sec:Ap} is not background.
The infinite product, the exact series and its two-sided bounds, the parity bound,
the near-critical asymptotic and the identification of $A(p)$ with a value of the
Koenigs linearising function of the logistic map are results of this thesis, placed
here rather than later because the quasi-stationary material of \cref{sec:qsd}
reduces to that single constant and because later chapters use it throughout. The
same applies to the closed form obtained in \cref{sec:absbirthdeath}. Everything
else in the chapter is standard, and is presented as such. Second, the argument is
in places heuristic --- continuum approximations to difference equations,
mean-field substitutions, asymptotic matching. Every such step is marked where it
occurs, and where a claim can be checked numerically it has been.
